%% bare_jrnl.tex
%% V1.3
%% 2007/01/11
%% by Michael Shell
%% see http://www.michaelshell.org/
%% for current contact information.
%%
%% This is a skeleton file demonstrating the use of IEEEtran.cls
%% (requires IEEEtran.cls version 1.7 or later) with an IEEE journal paper.
%%
%% Support sites:
%% http://www.michaelshell.org/tex/ieeetran/
%% http://www.ctan.org/tex-archive/macros/latex/contrib/IEEEtran/
%% and
%% http://www.ieee.org/



% *** Authors should verify (and, if needed, correct) their LaTeX system  ***
% *** with the testflow diagnostic prior to trusting their LaTeX platform ***
% *** with production work. IEEE's font choices can trigger bugs that do  ***
% *** not appear when using other class files.                            ***
% The testflow support page is at:
% http://www.michaelshell.org/tex/testflow/


%%*************************************************************************
%% Legal Notice:
%% This code is offered as-is without any warranty either expressed or
%% implied; without even the implied warranty of MERCHANTABILITY or
%% FITNESS FOR A PARTICULAR PURPOSE! 
%% User assumes all risk.
%% In no event shall IEEE or any contributor to this code be liable for
%% any damages or losses, including, but not limited to, incidental,
%% consequential, or any other damages, resulting from the use or misuse
%% of any information contained here.
%%
%% All comments are the opinions of their respective authors and are not
%% necessarily endorsed by the IEEE.
%%
%% This work is distributed under the LaTeX Project Public License (LPPL)
%% ( http://www.latex-project.org/ ) version 1.3, and may be freely used,
%% distributed and modified. A copy of the LPPL, version 1.3, is included
%% in the base LaTeX documentation of all distributions of LaTeX released
%% 2003/12/01 or later.
%% Retain all contribution notices and credits.
%% ** Modified files should be clearly indicated as such, including  **
%% ** renaming them and changing author support contact information. **
%%
%% File list of work: IEEEtran.cls, IEEEtran_HOWTO.pdf, bare_adv.tex,
%%                    bare_conf.tex, bare_jrnl.tex, bare_jrnl_compsoc.tex
%%*************************************************************************

% Note that the a4paper option is mainly intended so that authors in
% countries using A4 can easily print to A4 and see how their papers will
% look in print - the typesetting of the document will not typically be
% affected with changes in paper size (but the bottom and side margins will).
% Use the testflow package mentioned above to verify correct handling of
% both paper sizes by the user's LaTeX system.
%
% Also note that the "draftcls" or "draftclsnofoot", not "draft", option
% should be used if it is desired that the figures are to be displayed in
% draft mode.
%

\documentclass[journal]{IEEEtran}
%\documentclass[draftcls]{IEEEtran}


%\usepackage[english]{babel} 



%
% If IEEEtran.cls has not been installed into the LaTeX system files,
% manually specify the path to it like:
% \documentclass[journal]{../sty/IEEEtran}




% Some very useful LaTeX packages include:
% (uncomment the ones you want to load)


% *** MISC UTILITY PACKAGES ***
%
%\usepackage{ifpdf}
% Heiko Oberdiek's ifpdf.sty is very useful if you need conditional
% compilation based on whether the output is pdf or dvi.
% usage:
% \ifpdf
%   % pdf code
% \else
%   % dvi code
% \fi
% The latest version of ifpdf.sty can be obtained from:
% http://www.ctan.org/tex-archive/macros/latex/contrib/oberdiek/
% Also, note that IEEEtran.cls V1.7 and later provides a builtin
% \ifCLASSINFOpdf conditional that works the same way.
% When switching from latex to pdflatex and vice-versa, the compiler may
% have to be run twice to clear warning/error messages.






% *** CITATION PACKAGES ***
%
\usepackage{cite}
% cite.sty was written by Donald Arseneau
% V1.6 and later of IEEEtran pre-defines the format of the cite.sty package
% \cite{} output to follow that of IEEE. Loading the cite package will
% result in citation numbers being automatically sorted and properly
% "compressed/ranged". e.g., [1], [9], [2], [7], [5], [6] without using
% cite.sty will become [1], [2], [5]--[7], [9] using cite.sty. cite.sty's
% \cite will automatically add leading space, if needed. Use cite.sty's
% noadjust option (cite.sty V3.8 and later) if you want to turn this off.
% cite.sty is already installed on most LaTeX systems. Be sure and use
% version 4.0 (2003-05-27) and later if using hyperref.sty. cite.sty does
% not currently provide for hyperlinked citations.
% The latest version can be obtained at:
% http://www.ctan.org/tex-archive/macros/latex/contrib/cite/
% The documentation is contained in the cite.sty file itself.






% *** GRAPHICS RELATED PACKAGES ***
%
\ifCLASSINFOpdf
  \usepackage[pdftex]{graphicx}
  % declare the path(s) where your graphic files are
  \graphicspath{{img/}}
  % and their extensions so you won't have to specify these with
  % every instance of \includegraphics
  % \DeclareGraphicsExtensions{.pdf,.jpeg,.png}
\else
  % or other class option (dvipsone, dvipdf, if not using dvips). graphicx
  % will default to the driver specified in the system graphics.cfg if no
  % driver is specified.
  % \usepackage[dvips]{graphicx}
  % declare the path(s) where your graphic files are
  % \graphicspath{{../eps/}}
  % and their extensions so you won't have to specify these with
  % every instance of \includegraphics
  % \DeclareGraphicsExtensions{.eps}
\fi
% graphicx was written by David Carlisle and Sebastian Rahtz. It is
% required if you want graphics, photos, etc. graphicx.sty is already
% installed on most LaTeX systems. The latest version and documentation can
% be obtained at: 
% http://www.ctan.org/tex-archive/macros/latex/required/graphics/
% Another good source of documentation is "Using Imported Graphics in
% LaTeX2e" by Keith Reckdahl which can be found as epslatex.ps or
% epslatex.pdf at: http://www.ctan.org/tex-archive/info/
%
% latex, and pdflatex in dvi mode, support graphics in encapsulated
% postscript (.eps) format. pdflatex in pdf mode supports graphics
% in .pdf, .jpeg, .png and .mps (metapost) formats. Users should ensure
% that all non-photo figures use a vector format (.eps, .pdf, .mps) and
% not a bitmapped formats (.jpeg, .png). IEEE frowns on bitmapped formats
% which can result in "jaggedy"/blurry rendering of lines and letters as
% well as large increases in file sizes.
%
% You can find documentation about the pdfTeX application at:
% http://www.tug.org/applications/pdftex





% *** MATH PACKAGES ***
%
\usepackage[cmex10]{amsmath}
\usepackage{amssymb}
% A popular package from the American Mathematical Society that provides
% many useful and powerful commands for dealing with mathematics. If using
% it, be sure to load this package with the cmex10 option to ensure that
% only type 1 fonts will utilized at all point sizes. Without this option,
% it is possible that some math symbols, particularly those within
% footnotes, will be rendered in bitmap form which will result in a
% document that can not be IEEE Xplore compliant!
%
% Also, note that the amsmath package sets \interdisplaylinepenalty to 10000
% thus preventing page breaks from occurring within multiline equations. Use:
\interdisplaylinepenalty=2500
% after loading amsmath to restore such page breaks as IEEEtran.cls normally
% does. amsmath.sty is already installed on most LaTeX systems. The latest
% version and documentation can be obtained at:
% http://www.ctan.org/tex-archive/macros/latex/required/amslatex/math/





% *** SPECIALIZED LIST PACKAGES ***
%
%\usepackage{algorithmic}
% algorithmic.sty was written by Peter Williams and Rogerio Brito.
% This package provides an algorithmic environment fo describing algorithms.
% You can use the algorithmic environment in-text or within a figure
% environment to provide for a floating algorithm. Do NOT use the algorithm
% floating environment provided by algorithm.sty (by the same authors) or
% algorithm2e.sty (by Christophe Fiorio) as IEEE does not use dedicated
% algorithm float types and packages that provide these will not provide
% correct IEEE style captions. The latest version and documentation of
% algorithmic.sty can be obtained at:
% http://www.ctan.org/tex-archive/macros/latex/contrib/algorithms/
% There is also a support site at:
% http://algorithms.berlios.de/index.html
% Also of interest may be the (relatively newer and more customizable)
% algorithmicx.sty package by Szasz Janos:
% http://www.ctan.org/tex-archive/macros/latex/contrib/algorithmicx/




% *** ALIGNMENT PACKAGES ***
%
\usepackage{array}
% Frank Mittelbach's and David Carlisle's array.sty patches and improves
% the standard LaTeX2e array and tabular environments to provide better
% appearance and additional user controls. As the default LaTeX2e table
% generation code is lacking to the point of almost being broken with
% respect to the quality of the end results, all users are strongly
% advised to use an enhanced (at the very least that provided by array.sty)
% set of table tools. array.sty is already installed on most systems. The
% latest version and documentation can be obtained at:
% http://www.ctan.org/tex-archive/macros/latex/required/tools/


\usepackage{mdwmath}
\usepackage{mdwtab}
% Also highly recommended is Mark Wooding's extremely powerful MDW tools,
% especially mdwmath.sty and mdwtab.sty which are used to format equations
% and tables, respectively. The MDWtools set is already installed on most
% LaTeX systems. The lastest version and documentation is available at:
% http://www.ctan.org/tex-archive/macros/latex/contrib/mdwtools/


% IEEEtran contains the IEEEeqnarray family of commands that can be used to
% generate multiline equations as well as matrices, tables, etc., of high
% quality.


%\usepackage{eqparbox}
% Also of notable interest is Scott Pakin's eqparbox package for creating
% (automatically sized) equal width boxes - aka "natural width parboxes".
% Available at:
% http://www.ctan.org/tex-archive/macros/latex/contrib/eqparbox/





% *** SUBFIGURE PACKAGES ***
%\usepackage[tight,footnotesize]{subfigure}
% subfigure.sty was written by Steven Douglas Cochran. This package makes it
% easy to put subfigures in your figures. e.g., "Figure 1a and 1b". For IEEE
% work, it is a good idea to load it with the tight package option to reduce
% the amount of white space around the subfigures. subfigure.sty is already
% installed on most LaTeX systems. The latest version and documentation can
% be obtained at:
% http://www.ctan.org/tex-archive/obsolete/macros/latex/contrib/subfigure/
% subfigure.sty has been superceeded by subfig.sty.



%\usepackage[caption=false]{caption}
%\usepackage[font=footnotesize]{subfig}
% subfig.sty, also written by Steven Douglas Cochran, is the modern
% replacement for subfigure.sty. However, subfig.sty requires and
% automatically loads Axel Sommerfeldt's caption.sty which will override
% IEEEtran.cls handling of captions and this will result in nonIEEE style
% figure/table captions. To prevent this problem, be sure and preload
% caption.sty with its "caption=false" package option. This is will preserve
% IEEEtran.cls handing of captions. Version 1.3 (2005/06/28) and later 
% (recommended due to many improvements over 1.2) of subfig.sty supports
% the caption=false option directly:
%\usepackage[caption=false,font=footnotesize]{subfig}
%
% The latest version and documentation can be obtained at:
% http://www.ctan.org/tex-archive/macros/latex/contrib/subfig/
% The latest version and documentation of caption.sty can be obtained at:
% http://www.ctan.org/tex-archive/macros/latex/contrib/caption/




% *** FLOAT PACKAGES ***
%
%\usepackage{fixltx2e}
% fixltx2e, the successor to the earlier fix2col.sty, was written by
% Frank Mittelbach and David Carlisle. This package corrects a few problems
% in the LaTeX2e kernel, the most notable of which is that in current
% LaTeX2e releases, the ordering of single and double column floats is not
% guaranteed to be preserved. Thus, an unpatched LaTeX2e can allow a
% single column figure to be placed prior to an earlier double column
% figure. The latest version and documentation can be found at:
% http://www.ctan.org/tex-archive/macros/latex/base/



%\usepackage{stfloats}
% stfloats.sty was written by Sigitas Tolusis. This package gives LaTeX2e
% the ability to do double column floats at the bottom of the page as well
% as the top. (e.g., "\begin{figure*}[!b]" is not normally possible in
% LaTeX2e). It also provides a command:
%\fnbelowfloat
% to enable the placement of footnotes below bottom floats (the standard
% LaTeX2e kernel puts them above bottom floats). This is an invasive package
% which rewrites many portions of the LaTeX2e float routines. It may not work
% with other packages that modify the LaTeX2e float routines. The latest
% version and documentation can be obtained at:
% http://www.ctan.org/tex-archive/macros/latex/contrib/sttools/
% Documentation is contained in the stfloats.sty comments as well as in the
% presfull.pdf file. Do not use the stfloats baselinefloat ability as IEEE
% does not allow \baselineskip to stretch. Authors submitting work to the
% IEEE should note that IEEE rarely uses double column equations and
% that authors should try to avoid such use. Do not be tempted to use the
% cuted.sty or midfloat.sty packages (also by Sigitas Tolusis) as IEEE does
% not format its papers in such ways.


%\ifCLASSOPTIONcaptionsoff
%  \usepackage[nomarkers]{endfloat}
% \let\MYoriglatexcaption\caption
% \renewcommand{\caption}[2][\relax]{\MYoriglatexcaption[#2]{#2}}
%\fi
% endfloat.sty was written by James Darrell McCauley and Jeff Goldberg.
% This package may be useful when used in conjunction with IEEEtran.cls'
% captionsoff option. Some IEEE journals/societies require that submissions
% have lists of figures/tables at the end of the paper and that
% figures/tables without any captions are placed on a page by themselves at
% the end of the document. If needed, the draftcls IEEEtran class option or
% \CLASSINPUTbaselinestretch interface can be used to increase the line
% spacing as well. Be sure and use the nomarkers option of endfloat to
% prevent endfloat from "marking" where the figures would have been placed
% in the text. The two hack lines of code above are a slight modification of
% that suggested by in the endfloat docs (section 8.3.1) to ensure that
% the full captions always appear in the list of figures/tables - even if
% the user used the short optional argument of \caption[]{}.
% IEEE papers do not typically make use of \caption[]'s optional argument,
% so this should not be an issue. A similar trick can be used to disable
% captions of packages such as subfig.sty that lack options to turn off
% the subcaptions:
% For subfig.sty:
% \let\MYorigsubfloat\subfloat
% \renewcommand{\subfloat}[2][\relax]{\MYorigsubfloat[]{#2}}
% For subfigure.sty:
% \let\MYorigsubfigure\subfigure
% \renewcommand{\subfigure}[2][\relax]{\MYorigsubfigure[]{#2}}
% However, the above trick will not work if both optional arguments of
% the \subfloat/subfig command are used. Furthermore, there needs to be a
% description of each subfigure *somewhere* and endfloat does not add
% subfigure captions to its list of figures. Thus, the best approach is to
% avoid the use of subfigure captions (many IEEE journals avoid them anyway)
% and instead reference/explain all the subfigures within the main caption.
% The latest version of endfloat.sty and its documentation can obtained at:
% http://www.ctan.org/tex-archive/macros/latex/contrib/endfloat/
%
% The IEEEtran \ifCLASSOPTIONcaptionsoff conditional can also be used
% later in the document, say, to conditionally put the References on a 
% page by themselves.





% *** PDF, URL AND HYPERLINK PACKAGES ***
%
%\usepackage{url}
% url.sty was written by Donald Arseneau. It provides better support for
% handling and breaking URLs. url.sty is already installed on most LaTeX
% systems. The latest version can be obtained at:
% http://www.ctan.org/tex-archive/macros/latex/contrib/misc/
% Read the url.sty source comments for usage information. Basically,
% \url{my_url_here}.





% *** Do not adjust lengths that control margins, column widths, etc. ***
% *** Do not use packages that alter fonts (such as pslatex).         ***
% There should be no need to do such things with IEEEtran.cls V1.6 and later.
% (Unless specifically asked to do so by the journal or conference you plan
% to submit to, of course. )


% correct bad hyphenation here
\hyphenation{op-ti-cal net-works se-mi-con-duc-tor}


\renewcommand{\vec}[1]{\mathbf{#1}}
\newcommand{\mat}[1]{\mathbf{#1}}
\usepackage{color}
\newcommand{\hilight}[1]{\colorbox{yellow}{#1}}


\begin{document}

%\nocite{*}

%
% paper title
% can use linebreaks \\ within to get better formatting as desired
\title{Steady State Analysis of Nonlinear Dielectrics Filled Microwave Passive Devices}
%
%
% author names and IEEE memberships
% note positions of commas and nonbreaking spaces ( ~ ) LaTeX will not break
% a structure at a ~ so this keeps an author's name from being broken across
% two lines.
% use \thanks{} to gain access to the first footnote area
% a separate \thanks must be used for each paragraph as LaTeX2e's \thanks
% was not built to handle multiple paragraphs
%

\author{ \hilight{Order?} Laurent~Ntibarikure,~\IEEEmembership{Student Member,~IEEE,}
		Giuseppe~Pelosi,~\IEEEmembership{Fellow,~IEEE,}
		\hilight{ Lorenzo~Rossi},~%\IEEEmembership{Fellow,~IEEE,}
        and~Stefano~Selleri,~\IEEEmembership{Senior~Member,~IEEE}% <-this % stops a space
		\thanks{ L. Ntibarikure, G. Pelosi and S. Selleri are with
		Department of Electronics and Telecommunications,
		University of Florence, Via C. Lombroso 6/17 - 50134 - Florence, Italy } }%        
%\thanks{}%M. Shell is with the Department
%of Electrical and Computer Engineering, Georgia Institute of Technology, Atlanta,
%GA, 30332 USA e-mail: (see http://www.michaelshell.org/contact.html).}% <-this % stops a space
%\thanks{}%J. Doe and J. Doe are with Anonymous University.}% <-this % stops a space
%\thanks{}}%Manuscript received April 19, 2005; revised January 11, 2007.}}

% note the % following the last \IEEEmembership and also \thanks - 
% these prevent an unwanted space from occurring between the last author name
% and the end of the author line. i.e., if you had this:
% 
% \author{....lastname \thanks{...} \thanks{...} }
%                     ^------------^------------^----Do not want these spaces!
%
% a space would be appended to the last name and could cause every name on that
% line to be shifted left slightly. This is one of those "LaTeX things". For
% instance, "\textbf{A} \textbf{B}" will typeset as "A B" not "AB". To get
% "AB" then you have to do: "\textbf{A}\textbf{B}"
% \thanks is no different in this regard, so shield the last } of each \thanks
% that ends a line with a % and do not let a space in before the next \thanks.
% Spaces after \IEEEmembership other than the last one are OK (and needed) as
% you are supposed to have spaces between the names. For what it is worth,
% this is a minor point as most people would not even notice if the said evil
% space somehow managed to creep in.



% The paper headers
%\markboth{Journal of \LaTeX\ Class Files,~Vol.~6, No.~1, January~2007}%
\markboth{IEEE MICROWAVE AND WIRELESS COMPONENTS LETTERS, VOL. 20, NO. 6, JUNE 2010}%
{odd page header}
% The only time the second header will appear is for the odd numbered pages
% after the title page when using the twoside option.
% 
% *** Note that you probably will NOT want to include the author's ***
% *** name in the headers of peer review papers.                   ***
% You can use \ifCLASSOPTIONpeerreview for conditional compilation here if
% you desire.




% If you want to put a publisher's ID mark on the page you can do it like
% this:
%\IEEEpubid{0000--0000/00\$00.00~\copyright~2007 IEEE}
% Remember, if you use this you must call \IEEEpubidadjcol in the second
% column for its text to clear the IEEEpubid mark.



% use for special paper notices
%\IEEEspecialpapernotice{(Invited Paper)}




% make the title area
\maketitle


\begin{abstract}
%\boldmath
%The abstract goes here.
This letter presents the computation of high order harmonic components of the time-harmonic field exciting  waveguide structures partially filled with nonlinear dielectrics. The approach presented employs the Harmonic Balance method to perform a multi-harmonic Finite Element analysis, and the solution of the nonlinear system is handled with a Picard iteration. Since the nonlinear loop is particularly time demanding, a domain decomposition method is here proposed to speed-up the system solution, restricting the iteration loop to small sub-domain.

\end{abstract}
% IEEEtran.cls defaults to using nonbold math in the Abstract.
% This preserves the distinction between vectors and scalars. However,
% if the journal you are submitting to favors bold math in the abstract,
% then you can use LaTeX's standard command \boldmath at the very start
% of the abstract to achieve this. Many IEEE journals frown on math
% in the abstract anyway.

% Note that keywords are not normally used for peerreview papers.
\begin{IEEEkeywords}
% IEEEtran, journal, \LaTeX, paper, template.
harmonic balance, finite elements, domain decomposition, nonlinear dielectrics, waveguide filter
\end{IEEEkeywords}






% For peer review papers, you can put extra information on the cover
% page as needed:
% \ifCLASSOPTIONpeerreview
% \begin{center} \bfseries EDICS Category: 3-BBND \end{center}
% \fi
%
% For peerreview papers, this IEEEtran command inserts a page break and
% creates the second title. It will be ignored for other modes.
\IEEEpeerreviewmaketitle



\section{Introduction}
% The very first letter is a 2 line initial drop letter followed
% by the rest of the first word in caps.
% 
% form to use if the first word consists of a single letter:
% \IEEEPARstart{A}{demo} file is ....
% 
% form to use if you need the single drop letter followed by
% normal text (unknown if ever used by IEEE):
% \IEEEPARstart{A}{}demo file is ....
% 
% Some journals put the first two words in caps:
% \IEEEPARstart{T}{his demo} file is ....
% 
% Here we have the typical use of a "T" for an initial drop letter
% and "HIS" in caps to complete the first word.
%\IEEEPARstart{T}{his} demo file is intended to serve as a ``starter file''
%for IEEE journal papers produced under \LaTeX\ using
%IEEEtran.cls version 1.7 and later.
% You must have at least 2 lines in the paragraph with the drop letter
% (should never be an issue)
%I wish you the best of success.
%\hfill mds
 %\hfill January 11, 2007
% 
%Today's cutting edge microwave applications, due to the higher frequencies used and the growing power densities, have material nonlinear properties that become unveiled. 
\IEEEPARstart{T}{oday}'s microwave and millimeter wave passive devices design require modeling of material nonlinear properties to accurately predict the electromagnetic behavior, especially when high power densities are involved and dielectric materials are used. 
%Today's materials technologies and fabrication techniques have advanced considerably, increasing the need for accurate models predicting the electromagnetic behavior of designed devices. This is especially true when the device comprehend nonlinear materials, and steady state fields are searched for. 
A time-stepping approach to the solution of Maxwell's equations in such cases require considerable amount of simulation times. %\cite{deGersem2001}.
An alternative time-harmonic approach, based on the finite element (FE) solution of a partial differential equation (PDE) of time harmonic variable, provides the fields quantities at a given selected frequency. Unfortunately, a single frequency is not sufficient to capture all relevant phenomena, high order harmonics of the fields been generated from the interaction of the electromagnetic field with the nonlinear materials. Incorporating a larger set of harmonics in the FE simulation by using the Harmonic Balance Finite Element (HBFE) method \cite{Yamada1988,Yamada1989}, also known as or Multi-Harmonic Finite Element method \cite{Copeland2010, Bachinger2005, deGersem2001, Kolmbauer2009}, provides a fast and accurate solution to the problem, as long as the relevant harmonics required to affordably represent the field is not excessive. The nonlinearity of the PDE is tackled by an iterative solution of the so-built large-scale coupled linear system, updating the system at each step with its new linearized approach, converging within a prescribed tolerance to the sought solution. Several techniques are reported in literature, such as the Newton-Raphson algorithm. To improve the computational timings, an algebraic domain decomposition technique is applied as proposed in \cite{Guarnieri2010}, and significant improvements are achieved when the region occupied by nonlinear materials is rather small.

\par In this letter, we focus on the analysis of microwave devices filled with dielectrics presenting nonlinear scalar relative permittivity. 
%The presence of nonlinear materials in waveguide structures such as filters and ... renders finite elements (FE) simulations particularly time consuming, as an approach is used for the solution of the problem\cite{5582255}.
%A single frequency time-harmonic finite element model suffers from the inaccuracy due to the passive harmomic generation that occurs in nonlinear media, especially for high electromagnetic fields and highly nonlinear media. To overcome this problem, an Harmonic Balance (HB) approach has to be used. 
%The harmonic balance in finite element modeling was first introduced by Yamada et Al. in magnetostatics, for the computation of the harmonics generated within ferromagnetic media which presents a magnetic flux saturation polynomial behavior \cite{92182}. We apply the technique to an electromagnetic wave problem, obtaining the coupling coefficents between selected harmonics.
In section \ref{sec:formulation}, we present the Harmonic Balance Finite Element method formulation for the solution of the vector Helmholtz equation for the electric field, within a domain comprising nonlinear dielectrics. In section \ref{sec:results}, simulation results of a broadband E-plane bandpass filter filled with Kerr-type material are shown, and performances of the algorithm are discussed.

%%%%%%%%%%%%%%%%%%%%%%%%%%%%%%%%%%%%%%%%%%%%%%%%%%%%%%%%%%%%%%%%%%%%%%%%%%%%%%%%%%%%%%%%%%%%%%%%%%%%%%%%%%%%%%%%%%%%%%%%%%
\section{Formulation} \label{sec:formulation}


\begin{figure}[!t]
\centering
\includegraphics[width=0.6\columnwidth]{geometry}
\caption{Sketch of the domain $\Omega = \Omega_1 \cup \Omega_2$. \hilight{increase font size}}
\label{fig:geometry}
\end{figure}

\par The formulation given in this section holds for the case of a generic electromagnetic problem defined on a closed domain as depicted in Fig. \ref{fig:geometry}. The domain $\Omega$, with boundary $\Gamma_0$, can be either defined on $\mathbb R^2$ or $\mathbb R^3$. For the sake of simplicity, the examples presented in the next section are analyzed in $\mathbb R^2$, allowing for straightforward implementation. $\Omega$ is divided into a large part, $\Omega_1$, with linear scalar relative permittivity $\bar{\varepsilon}_r(\vec{r})$ and a small part, $\Omega_2$, with a nonlinear relative permittivity $\tilde{\varepsilon}_r(\vec{E}(\vec{r}))$.
%Consider a general electromagnetic problem defined on a domain $\Omega$ of the boundary $\Gamma_0$. $\Omega$ is chosen to be $\mathbb R^2$, but the extension to the $\mathbb R^3$ case is straightforward. The scalar relative permittivity of the widest part of the domain $\Omega_1$ is linear and equal to $\bar{\varepsilon}_r(\vec{r})$, while a small part of the domain $\Omega_2$ (not necessarily simply connected), exhibits a nonlinear permittivity $\tilde{\varepsilon}_r(\vec{E}(\vec{r}))$. 
The scalar permittivity in $\Omega$ can be expressed as 
\begin{equation}
\varepsilon_r(\vec{E}(\vec{r}),\vec{r}) = 
\begin{cases}
\bar{\varepsilon}_r(\vec{r}), & \vec{r} \in \Omega_1 \\
\tilde{\varepsilon}_r(\vec{E}(\vec{r}), \vec{r}) = \bar{\varepsilon}_r(\vec{r}) + \mathcal{L}(\vec{E}(\vec{r})), & \vec{r} \in \Omega_2
\end{cases}
\end{equation} where $\vec{E}(\vec{r})$ is the electric field pointed by $\vec{r}$ and $\mathcal{L}(\cdot)$ denotes a generic scalar operator that describes the nonlinear behavior. The electric field satisfies, within $\Omega$, the vector Helmholtz equation
\begin{equation}
\label{eq:helmholtz}
\nabla \times \left ( \frac{1}{\mu_r} \nabla \times \vec{E}(\vec{r}) \right )  - {k_0}^2 \varepsilon_r \left(\vec{E}(\vec{r}), \vec{r} \right) \vec{E}(\vec{r}) = 0
\end{equation} with suitable boundary conditions on $\Gamma_0$. For the example shown in section \ref{sec:results}, $\Gamma_0$ is considered partly perfect electric conductor, partly a modal expansion of the field at the wave ports as developed in \cite{Pelosi2009}. 

\par We proceed with the standard FE Galerkin framework: the domain $\Omega$ is discretized into a set of non overlapping finite elements and the electric field approximated as a linear combination of vector basis functions $\boldsymbol{\alpha}_n$ 
\begin{equation}
\vec{E}(\vec{r})  = \sum_{n=1}^N E_n \boldsymbol{\alpha}_n(\vec{r})
\end{equation} % and the weighting functions $\vec{W}_m$ are the same as the shape functions $\boldsymbol{\alpha}_n$. As a result of testing of \eqref{eq:helmholtz} with the weighting functions, we obtain the following weak form
Then by testing \eqref{eq:helmholtz} over $\Omega$ with the weighting functions $\vec{W}_m$, chosen the same as the shape functions $\boldsymbol{\alpha}_n$, we obtain the following weak form:
\begin{eqnarray}\label{eq:fem}
%\int_\Omega \vec{W}_m (\vec{r}) \cdot \left \{ \nabla \times \left ( \frac{1}{\mu_r} \nabla \times \left ( \sum_n E_n \boldsymbol{\alpha}_n(\vec{r})\right ) \right ) \right. \nonumber \\
%\left. - {k_0}^2 \varepsilon_r (\vec{E}(\vec{r})) \left ( \sum_n E_n \boldsymbol{\alpha}_n(\vec{r}) \right ) \right \} d\Omega = 0
& & \hspace{-1.5cm} \int_\Omega \nabla \times \vec{W}_m (\vec{r}) \cdot \frac{1}{\mu_r} \nabla \times \vec{E}(\vec{r}) \ d\Omega \nonumber \\
& & \hspace{-.75cm} - k_0^2 \int_\Omega  \vec{W}_m (\vec{r}) \cdot \varepsilon_r (\vec{E}(\vec{r}), \vec{r}) \vec{E}(\vec{r}) \ d\Omega \nonumber \\
& & + \oint_{\Gamma_0} \vec{W}_m (\vec{r}) \cdot \hat{n} \times \frac{1}{\mu_r} \nabla \times \vec{E}(\vec{r}) \ d\Gamma = 0
%\left [ \sum_n E_n \boldsymbol{\alpha}_n(\vec{r}) \right ]
\end{eqnarray} 

%%%%% HBFE  %%%%%

\par We introduce the multi-harmonic dependence of the HBFE method, by approximating the electric field as a truncated Fourier series \cite{Bachinger2005}
\begin{equation} \label{eq:field}
%\vec{E} = \vec{E}_{1s} \sin(\omega t) + \vec{E}_{1c} \cos(\omega t) + \vec{E}_{3s} \sin(3 \omega t) + \vec{E}_{3c} \cos(3 \omega t)
\vec{\hat{E}}(\vec{r}, t) = \sum^M_{i=1} \vec{E}^{is}(\vec{r}) \sin(i \omega t) + \vec{E}^{ic}(\vec{r}) \cos(i \omega t)
\end{equation} where $\omega$ is the fundamental angular frequency. $\vec{E}_{i}^{s}$ and $\vec{E}_{i}^{c}$ are the harmonic complex amplitudes which expand the finite element coefficients $E_n$, which must individually satisfy \eqref{eq:helmholtz}. 
%For each of the harmonic components retained in the truncated series, the field must satisfy \eqref{eq:helmholtz}.
The scalar relative permittivity can also be approximated as a truncated Fourier series 
\begin{eqnarray} \label{eq:eps}
\hat{{\varepsilon}}_r(\vec{E}(\vec{r}), \vec{r}, t) = {\varepsilon}_{r_{0}} + \sum^K_{k=1} {\varepsilon}_{r}^{ks} \sin(k \omega t) + {\varepsilon}_{r}^{kc} \cos(k \omega t)
\end{eqnarray} and the coefficients of the expansion are given by
\begin{eqnarray} \label{eq:epsCoeffs}
{\varepsilon}_{r_{0}} & = & \frac{\omega}{2\pi}\int_0^{\frac{2\pi}{\omega}} {\varepsilon}_r(\vec{E}(\vec{r}), \vec{r}) \ dt \\
{\varepsilon}_{r}^{ks} & = & \frac{\omega}{\pi}\int_0^{\frac{2\pi}{\omega}} {\varepsilon}_r(\vec{E}(\vec{r}), \vec{r}) \sin(k \omega t) \ dt \\
{\varepsilon}_{r}^{kc} & = & \frac{\omega}{\pi}\int_0^{\frac{2\pi}{\omega}} {\varepsilon}_r(\vec{E}(\vec{r}), \vec{r}) \cos(k \omega t) \ dt
\end{eqnarray} with $k = 1 \dots K$. The orders $M$ and $K$ of approximation of the electric field and the permittivity, respectively, depend on the specific device analyzed and on the intensity of the incident field. These can be set by choosing an appropriate convergence criterion of the power associated to the harmonics of lower order as the number of multi-harmonic coefficients is increased.

\par We proceed substituting $\vec{E}(\vec{r})$ and ${\varepsilon}_r(\vec{r})$ by their respective harmonic approximations $\vec{\hat{E}}(\vec{r}, t)$ and $\hat{{\varepsilon}}_r(\vec{r}, t)$ in \eqref{eq:fem}, then testing the equation with $\sin(i \omega t)$ and $\cos(i \omega t)$, $i = 1 \dots M$, and integrating over the time period of the fundamental $\left [ 0,\frac{2\pi}{\omega} \right ]$ gives a large-scale linear system of the form:
\begin{equation}
\scriptscriptstyle
\left[ \begin{array}{cccc}
a_{11}^{1s1s} & a_{11}^{1s1c} & \ldots &  a_{1N}^{1sMc} \\
a_{11}^{1c1s} & a_{11}^{1c1c} & \ldots &  a_{1N}^{1cMc} \\
\vdots & \vdots & \ddots & \vdots \\
a_{N1}^{Mc1s} & a_{N1}^{Mc1c} & \ldots &  a_{NN}^{McMc}
\end{array}\right] 
\left[ \begin{array}{cccc}
E_{1}^{1s} \\
E_{1}^{1c} \\
\vdots \\
E_{N}^{Mc}
\end{array}\right] =
\left[ \begin{array}{cccc}
b_{1}^{1s} \\
b_{1}^{1c} \\
\vdots \\
b_{N}^{Mc}
\end{array}\right] 
\end{equation}
The product that occurs in the second integral of \ref{eq:fem}, between the approximated field and permittivity, induces coupling between the harmonic components. The order of the resulting HBFE system is of $N M$. In practical analyses, this number can become considerable, and as introduced previously, the solutions have to be computed iteratively, building at each step a new linearized system, until convergence is achieved. The chosen nonlinear solution is based on a Picard iteration: the permittivity is assumed initially to be the one in the absence of electric field: 
\begin{equation} \label{eq:firstLinearSystem}
[A_{([E]^0)}] [E]^1 = [b_{([E]^0)}]
\end{equation} where $[E]^0$ is a null vector of unknowns. The solution of the first linear system of \eqref{eq:firstLinearSystem} provides expansion coefficients that will be used to compute a new value of the permittivity:
\begin{equation} \label{eq:firstLinearSystem}
[A_{([E]^{q-1})}] [E]^{q} = [b_{([E]^{q-1})}]
\end{equation}
 The process is repeated until the updated solution do not differ from the previous one, in the sense of the euclidean norm, more than a prescribed value.
\par The Picard iteration is not always stably convergent, especially for highly nonlinear materials and high impressed fields. Improvements can be achieved updating the system using fields solutions given by a proper linear combination of previous solutions \cite{Silvester1997}.

\par To enhance the computational efficiency, a substructuring domain decomposition approach is used \cite{Toselli2005, Guarnieri2009, Guarnieri2010, Saad2000}. The multiharmonic unknown vector $[E]$ defined for the HBFE method over $\Omega$ can be split into two vectors $[E^{(1)}]$ and $[E^{(2)}]$
containing the unknowns belonging to interior points in $\Omega_1$ and $\Omega_2$, respectively. Unknowns belonging to the boundaries are placed in a third vector $[E^c]$. The HBFE system can be recast in the form:
\begin{equation}\label{eq:DD}
\scriptstyle
\left [ \begin{array}{cccccc}
\scriptstyle [M^{(1)}] & \scriptstyle 0  &  \scriptstyle[P^{(1)}] \\
\scriptstyle 0  &  \scriptstyle[M_{([E]^{q-1})}^{(2)}] & \scriptstyle[P_{([E]^{q-1})}^{(2)}] \\
\scriptstyle{[Q^{(1)}]} & \scriptstyle{[Q_{([E]^{q-1})}^{(2)}]} & \scriptstyle{[C_{([E]^{q-1})}]}
\end{array} \right ]
\left [ \begin{array}{cc}
\scriptstyle{[E^{(1)}]^{q}} \\
\scriptstyle{[E^{(2)}]^{q}} \\
\scriptstyle{[E^{(c)}]^{q}}
\end{array} \right ] = \left [ \begin{array}{cc}
\scriptstyle{[b^{(1)}]} \\
\scriptstyle{[b_{([E]^{q-1})}^{(2)}]} \\
\scriptstyle{[b_{([E]^{q-1})}^{(c)}]}
\end{array} \right ] 
\end{equation}
where $[M^{(1)}]$ contains the FE coefficients of the linear system related to the unknowns $[E^{(1)}]$ fully internal for $\Omega_1$, thus linear, and $[M^{(2)}]$ contains coefficients for the nonlinear sub-domain $\Omega_2$. ${[P^{(\cdot)}]}$ and ${[Q^{(\cdot)}]}$ represent couplings between the interior unknowns and the ones on the boundary $\Gamma_{12}$ collected in $[E^{(c)}]$. Using a Gauss elimination-type strategy as the Schur complement $[S_{([E]^{q-1})}]$, the boundary unknowns $[E^{(c)}]^{q}$ of \eqref{eq:DD} are computed by solving the reduced system
\begin{equation}
[S_{([E]^{q-1})}] [E^{(c)}]^{q} = [d]^{q-1} 
\end{equation} where
\begin{eqnarray}
\label{eq:shurS}
& & \hspace{-1cm} [S_{([E]^{q-1})}] = [C_{([E]^{q-1})}] - [Q^{(1)}][M^{(1)}]^{-1} [P^{(1)}] \nonumber \\
& & \hspace{1cm}- [Q_{([E]^{q-1})}^{(2)}][M_{([E]^{q-1})}^{(1)}]^{-1} [P_{([E]^{q-1})}^{(1)}]
\end{eqnarray} and
\begin{eqnarray}
\label{eq:shurb}
& & \hspace{-1cm} [d]^k = [b_{([E]^{q-1})}^{(c)}] - [Q^{(1)}][M^{(1)}]^{-1} [P^{(1)}] \nonumber \\
& & \hspace{1cm}- [Q_{([E]^{q-1})}^{(2)}][M_{([E]^{q-1})}^{(1)}]^{-1} [P_{([E]^{q-1})}^{(1)}]
\end{eqnarray}
The difference operators in \eqref{eq:shurS} and \eqref{eq:shurb} are to be intended as the assembling of local matrices in a global one as performed by global FE matrix assembly \cite{Guarnieri2009}. Once $[E^{(c)}]^{q}$ is known, local solutions are given by
\begin{eqnarray}
[E^{(1)}]^{q-1} & = & [M^{(1)}]^{-1} \left( [b^{(1)}] - [P^{(1)}] {[E^{(c)}]^{q}} \right) \\
{[E^{(2)}]}^{q-1} & = & [M_{([E]^{q-1})}^{(2)}]^{-1} \left( [ b_{([E]^{q-1})}^{(2)} ] \right. \nonumber \\
& & \hspace{1cm} \left. - [P_{([E]^{q-1})}^{(2)}] [E^{(c)}]^{q} \right)
\end{eqnarray}
In order to solve the HBFE system and compute the steady state field, every single submatrix is assembled at first sight, then, within the iteration loop, only the submatrices related to $\Omega_2$ are to be updated. The efficiency of this DD technique relies on the fragmentation of the matrices to invert. When the number of coefficients related to nonlinear media is small, as usually happens for microwave devices where the domain of nonlinear media has reduced extension, noticeable improvements are achieved.


\section{Numerical Results} \label{sec:results}

\par In this section, we analyze the microwave filter sketched in Fig. \ref{fig:geometry}. The dielectric material
\cite{Yatsyk2006} (Yatsyk2006)

\begin{figure}[!t]
\centering
\includegraphics[width=\columnwidth]{bilat}
\caption{Millimeter wave filter. \hilight{increase font size}}
\label{fig:geometry}
\end{figure}

%\_subsection{Subsection Heading Here}
%Subsection text here.

% needed in second column of first page if using \IEEEpubid
%\IEEEpubidadjcol

%\_subsubsection{Subsubsection Heading Here}
%Subsubsection text here.


% An example of a floating figure using the graphicx package.
% Note that \label must occur AFTER (or within) \caption.
% For figures, \caption should occur after the \includegraphics.
% Note that IEEEtran v1.7 and later has special internal code that
% is designed to preserve the operation of \label within \caption
% even when the captionsoff option is in effect. However, because
% of issues like this, it may be the safest practice to put all your
% \label just after \caption rather than within \caption{}.
%
% Reminder: the "draftcls" or "draftclsnofoot", not "draft", class
% option should be used if it is desired that the figures are to be
% displayed while in draft mode.
%
%\begin{figure}[!t]
%\centering
%\includegraphics[width=2.5in]{myfigure}
% where an .eps filename suffix will be assumed under latex, 
% and a .pdf suffix will be assumed for pdflatex; or what has been declared
% via \DeclareGraphicsExtensions.
%\caption{Simulation Results}
%\label{fig_sim}
%\end{figure}

% Note that IEEE typically puts floats only at the top, even when this
% results in a large percentage of a column being occupied by floats.


% An example of a double column floating figure using two subfigures.
% (The subfig.sty package must be loaded for this to work.)
% The subfigure \label commands are set within each subfloat command, the
% \label for the overall figure must come after \caption.
% \hfil must be used as a separator to get equal spacing.
% The subfigure.sty package works much the same way, except \subfigure is
% used instead of \subfloat.
%
%\begin{figure*}[!t]
%\centerline{\subfloat[Case I]\includegraphics[width=2.5in]{subfigcase1}%
%\label{fig_first_case}}
%\hfil
%\subfloat[Case II]{\includegraphics[width=2.5in]{subfigcase2}%
%\label{fig_second_case}}}
%\caption{Simulation results}
%\label{fig_sim}
%\end{figure*}
%
% Note that often IEEE papers with subfigures do not employ subfigure
% captions (using the optional argument to \subfloat), but instead will
% reference/describe all of them (a), (b), etc., within the main caption.


% An example of a floating table. Note that, for IEEE style tables, the 
% \caption command should come BEFORE the table. Table text will default to
% \footnotesize as IEEE normally uses this smaller font for tables.
% The \label must come after \caption as always.
%
%\begin{table}[!t]
%% increase table row spacing, adjust to taste
%\renewcommand{\arraystretch}{1.3}
% if using array.sty, it might be a good idea to tweak the value of
% \extrarowheight as needed to properly center the text within the cells
%\caption{An Example of a Table}
%\label{table_example}
%\centering
%% Some packages, such as MDW tools, offer better commands for making tables
%% than the plain LaTeX2e tabular which is used here.
%\begin{tabular}{|c||c|}
%\hline
%One & Two\\
%\hline
%Three & Four\\
%\hline
%\end{tabular}
%\end{table}


% Note that IEEE does not put floats in the very first column - or typically
% anywhere on the first page for that matter. Also, in-text middle ("here")
% positioning is not used. Most IEEE journals use top floats exclusively.
% Note that, LaTeX2e, unlike IEEE journals, places footnotes above bottom
% floats. This can be corrected via the \fnbelowfloat command of the
% stfloats package.



\section{Conclusion}
The conclusion goes here.





% if have a single appendix:
%\appendix[Proof of the Zonklar Equations]
% or
%\appendix  % for no appendix heading
% do not use \section anymore after \appendix, only \section*
% is possibly needed

% use appendices with more than one appendix
% then use \section to start each appendix
% you must declare a \section before using any
% \subsection or using \label (\appendices by itself
% starts a section numbered zero.)
%


%\appendices
%\_section{Proof of the First Zonklar Equation}
%Appendix one text goes here.

% you can choose not to have a title for an appendix
% if you want by leaving the argument blank
%\_section{}
%Appendix two text goes here.


% use section* for acknowledgement
%\_section*{Acknowledgment}


%The authors would like to thank...


% Can use something like this to put references on a page
% by themselves when using endfloat and the captionsoff option.
\ifCLASSOPTIONcaptionsoff
  \newpage
\fi



% trigger a \newpage just before the given reference
% number - used to balance the columns on the last page
% adjust value as needed - may need to be readjusted if
% the document is modified later
%\IEEEtriggeratref{8}
% The "triggered" command can be changed if desired:
%\IEEEtriggercmd{\enlargethispage{-5in}}

% references section

% can use a bibliography generated by BibTeX as a .bbl file
% BibTeX documentation can be easily obtained at:
% http://www.ctan.org/tex-archive/biblio/bibtex/contrib/doc/
% The IEEEtran BibTeX style support page is at:
% http://www.michaelshell.org/tex/ieeetran/bibtex/
%\bibliographystyle{IEEEtran}
% argument is your BibTeX string definitions and bibliography database(s)
%\bibliography{IEEEabrv,../bib/paper}
%
% <OR> manually copy in the resultant .bbl file
% set second argument of \begin to the number of references
% (used to reserve space for the reference number labels box)
%\begin{thebibliography}{1}
%
%\bibitem{IEEEhowto:kopka}
%H.~Kopka and P.~W. Daly, \emph{A Guide to \LaTeX}, 3rd~ed.\hskip 1em plus
%  0.5em minus 0.4em\relax Harlow, England: Addison-Wesley, 1999.
%
%\end{thebibliography}
\bibliographystyle{IEEEtran}	% (uses file "plain.bst")
\bibliography{HBFEDD}


% biography section
% 
% If you have an EPS/PDF photo (graphicx package needed) extra braces are
% needed around the contents of the optional argument to biography to prevent
% the LaTeX parser from getting confused when it sees the complicated
% \includegraphics command within an optional argument. (You could create
% your own custom macro containing the \includegraphics command to make things
% simpler here.)
%\begin{biography}[{\includegraphics[width=1in,height=1.25in,clip,keepaspectratio]{mshell}}]{Michael Shell}
% or if you just want to reserve a space for a photo:

%\begin{IEEEbiography}{Michael Shell}
%Biography text here.
%\end{IEEEbiography}

% if you will not have a photo at all:
%\begin{IEEEbiographynophoto}{John Doe}
%Biography text here.
%\end{IEEEbiographynophoto}

% insert where needed to balance the two columns on the last page with
% biographies
%\newpage

%\begin{IEEEbiographynophoto}{Jane Doe}
%Biography text here.
%\end{IEEEbiographynophoto}

% You can push biographies down or up by placing
% a \vfill before or after them. The appropriate
% use of \vfill depends on what kind of text is
% on the last page and whether or not the columns
% are being equalized.

%\vfill

% Can be used to pull up biographies so that the bottom of the last one
% is flush with the other column.
%\enlargethispage{-5in}



% that's all folks
\end{document}


